% Lane WA — theory. All TikZ libraries used here are loaded in main.tex's
% preamble (shapes.geometric, arrows.meta, positioning, calc, fit, backgrounds,
% decorations.pathmorphing, decorations.pathreplacing, matrix); no additions needed.
\section{Theory}\label{sec:theory}

This section develops the four ideas the compiler is built on: tile-native spatial
compute (Section~\ref{sec:theory} kernel model), operator fusion, NoC-aware sharding,
and the numerical-validation theory behind the correctness bar. Every constant, byte
count, and measured value quoted here is taken from the repository
(\texttt{lib/tensixflow/Scheduler/L1Model.h}, \texttt{docs/performance\_model.md},
\texttt{reports/rmsnorm\_matmul\_plan.json}, and the checked-in test evidence);
nothing is quoted from outside it.

\subsection{Tile-native spatial compute}

The modeled device is a single Wormhole n150-class part: an 8x8 grid of 64
Tensix-style cores. Each core is a tile-native engine: it does not operate on raw
tensors but on 32x32 element tiles, with bf16 data and fp32 accumulation (modeled
math fidelity HiFi4). Every core owns a private slice of L1 SRAM; tt-metal-style
memory addressing is \texttt{(x, y, local L1 address)}, so any core can name any
other core's L1. Cores are connected by a Network-on-Chip (NoC) supporting
asynchronous reads, writes, and multicast; DRAM sits off-chip behind the NoC.

A program placed on a core is decomposed into three kernel roles, matching the
TT-Metalium convention:

\begin{itemize}[nosep]
  \item \textbf{Reader} (data movement): pulls input tiles from DRAM (or from a
        peer core's L1) over the NoC and pushes them into input circular buffers.
  \item \textbf{Compute}: pops tiles from the input buffers, runs the tile math
        (rmsnorm, matmul, gelu), and pushes result tiles into the output buffer.
  \item \textbf{Writer} (data movement): drains the output circular buffer and
        writes the tiles back over the NoC to DRAM.
\end{itemize}

The three roles never touch each other directly. They communicate exclusively
through \emph{circular buffers}: on-chip producer/consumer queues of tiles with a
fixed depth. The push side reserves space
(\texttt{cb\_reserve\_back}/\texttt{cb\_push\_back}) and blocks when the queue is
full; the pop side waits for data
(\texttt{cb\_wait\_front}/\texttt{cb\_pop\_front}) and blocks when the queue is
empty. Depth is the pipelining knob: it controls how far ahead the reader may run
of the compute (and the compute of the writer) before stalling.

Figure~\ref{fig:theory-core-memory} shows the memory view of one modeled core. The
per-core L1 budget is the single shared constant set of
\texttt{lib/tensixflow/Scheduler/L1Model.h}: 1\,MiB total
(\texttt{kL1TotalBytes} $= 1048576$) minus a 64\,KiB reserve for code and stack
(\texttt{kL1ReserveBytes} $= 65536$) leaves a 960\,KiB circular-buffer + shard
budget (\texttt{kCBShardBudget} $= 983040$). One 32x32 bf16 tile costs
$32 \times 32 \times 2 = 2048$ bytes (\texttt{kTileBytes}). The default
three-buffer pipeline ($c_0$ = in0, $c_1$ = in1, $c_{16}$ = out, the tt-metal
index convention) then runs at depth
$\lfloor 983040 / (3 \times 2048) \rfloor = 160$ tiles per buffer, consuming
$3 \times 160 \times 2048 = 983040$ bytes --- exactly the budget, i.e. occupancy
1.0.

\begin{figure}[t]
  \centering
  \begin{tikzpicture}[
      font=\small,
      tile/.style={draw=tensixteal, thick, fill=tensixlight,
                   minimum width=1.35cm, minimum height=4.5mm, inner sep=0pt},
      cblab/.style={font=\footnotesize\bfseries, text=tensixnavy},
      cbsub/.style={font=\scriptsize, text=tensixslate},
      note/.style={font=\footnotesize, text=tensixnavy, align=left},
    ]
    % ---- Left: per-core L1 bar, height proportional to bytes ----------------
    \def\barw{2.4} \def\barh{6.4}   % 6.4 cm = 1024 KiB; reserve = 64 KiB = 0.4 cm
    \coordinate (base) at (0,0);
    \fill[tensixlight] (base) rectangle ++(\barw,\barh);
    \fill[tensixslate!35] ($(base)+(0,\barh-0.4)$) rectangle ++(\barw,0.4);
    \draw[tensixnavy, very thick] (base) rectangle ++(\barw,\barh);
    \draw[tensixnavy, thick] ($(base)+(0,\barh-0.4)$) -- ++(\barw,0);
    \node[font=\tiny, text=tensixnavy] at ($(base)+(0.5*\barw,\barh-0.2)$) {reserve};
    \node[font=\scriptsize, rotate=90, text=tensixnavy]
      at ($(base)+(0.5*\barw,0.5*\barh-0.2)$) {CB + shard budget (960\,KiB)};

    % Brace left: total budget
    \draw[decorate, decoration={brace, amplitude=5pt}, thick, tensixslate]
      ($(base)+(0,\barh)$) -- (base)
      node[midway, left=8pt, align=right, font=\footnotesize, text=tensixnavy]
        {L1 total\\$=$ 1\,MiB\\$=$ 1\,048\,576\,B};
    % Braces right: reserve and CB/shard budget
    \draw[decorate, decoration={brace, amplitude=4pt}, thick, tensixslate]
      ($(base)+(\barw,\barh)$) -- ++(0,-0.4)
      node[midway, right=7pt, font=\scriptsize, text=tensixnavy, align=left]
        {64\,KiB code/stack\\reserve};
    \draw[decorate, decoration={brace, amplitude=4pt}, thick, tensixslate]
      ($(base)+(\barw,\barh-0.4)$) -- ($(base)+(\barw,0)$)
      node[midway, right=7pt, font=\scriptsize, text=tensixnavy, align=left]
        {960\,KiB\\circular-buffer\\+ shard budget};

    % ---- Right: the three circular buffers as tile queues --------------------
    \foreach \x/\idx/\role in {5.6/c_0/in0 (reader fills),
                               7.6/c_1/in1 (reader fills),
                               9.6/c_{16}/out (writer drains)} {
      \node[cblab] at (\x, 3.55) {$\idx$};
      \node[cbsub, align=center] at (\x, 3.15) {\role};
      \foreach \dy in {0, 0.5, 1.0} {
        \node[tile] at (\x, 0.35+\dy) {};
      }
      \node[font=\footnotesize, text=tensixslate] at (\x, 1.85) {$\vdots$};
      \node[tile] at (\x, 2.25) {};
      \draw[decorate, decoration={brace, amplitude=3pt}, thin, tensixslate]
        (\x+0.78, 0.12) -- (\x+0.78, 2.48)
        node[midway, right=4pt, font=\tiny, text=tensixslate] {depth 160};
    }
    % compute flow between buffers
    \draw[-Stealth, tensixteal, thick] (6.0, 2.75) .. controls (6.8, 3.1) .. (7.2, 2.75);
    \node[font=\tiny, text=tensixteal] at (6.6, 3.0) {compute};
    \draw[-Stealth, tensixteal, thick] (8.0, 2.75) .. controls (8.8, 3.1) .. (9.2, 2.75);

    % ---- Bottom: the arithmetic ------------------------------------------------
    \node[note] at (6.9, -0.45)
      {tile $= 32 \times 32 \times 2 = 2048$\,B\quad
       depth $= \lfloor 983040 / (3 \cdot 2048) \rfloor = 160$\quad
       $3 \times 160 \times 2048 = 983040$\,B $=$ budget (occupancy 1.0)};

    % background grouping for the CB region
    \begin{scope}[on background layer]
      \node[fit={($(5.6,0)+( -0.85,-0.05)$) ($(9.6,3.9)+(0.85,0.05)$)},
            draw=tensixslate, thin, densely dotted, fill=tensixlight!40,
            inner sep=2pt] {};
    \end{scope}
  \end{tikzpicture}
  \caption{Memory view of one modeled core. Left: the per-core L1 budget of
    \texttt{L1Model.h} --- 1\,MiB total minus a 64\,KiB code/stack reserve leaves
    the 960\,KiB circular-buffer + shard budget. Right: the default three-buffer
    pipeline ($c_0$ in0, $c_1$ in1, $c_{16}$ out) drawn as on-chip
    producer/consumer tile queues at depth 160; the three buffers consume exactly
    the budget ($983040$\,B, occupancy 1.0). Tile size 2048\,B is bf16 32x32.}
  \label{fig:theory-core-memory}
\end{figure}

\subsection{Operator fusion}

Without fusion, every operator boundary round-trips DRAM: RMSNorm reads the
activation from DRAM and writes the normalized intermediate back out, and the
following MatMul re-reads that same intermediate from DRAM (plus the weight) before
writing its output. Fusing the pair into one kernel keeps the intermediate resident
in L1 between the two stages --- the rmsnorm output never leaves the core that
produced it --- and erases both the intermediate DRAM write and the intermediate
DRAM re-read.

For the running example \texttt{rmsnorm\_matmul} (activation
$[1, 2048, 4096]$ bf16, weight $[4096, 12288]$ bf16) the intermediate has the same
shape as the activation, so it costs $1 \times 2048 \times 4096 \times 2 =
16\,777\,216$ bytes. That is exactly the figure the planner records in the plan
JSON as \texttt{estimated\_dram\_reads\_saved} (and
\texttt{estimated\_intermediate\_writes\_saved}): \textbf{16\,777\,216 bytes} each,
with the eliminated tensor counted on the fused op via the
\texttt{tensixflow.eliminated\_bytes} / \texttt{tensixflow.eliminated\_tensors}
annotations. Figure~\ref{fig:theory-fusion-dataflow} contrasts the two dataflows.

The fusion pass goes one step further on the full MLP chain
(\texttt{examples/mlp\_block.mlir}):
\[
  \text{rmsnorm} \to \text{matmul}(W_1) \to \text{add}(b_1) \to \text{gelu}
  \to \text{matmul}(W_2) \to \text{add}(\text{residual})
\]
is collapsed, link by link with a single-use guard at every intermediate, into a
single \texttt{tensixflow.fused\_mlp} kernel --- one kernel launch instead of six
operators, with all intermediates circulating in L1/registers rather than DRAM.%
\footnote{Grounding reference (inline per the citation-free contract):
  arXiv:2606.09879, \emph{Operator Fusion for LLM Inference on the Tensix
  Architecture} --- fuses RMSNorm with matmul in on-chip SRAM to cut DRAM
  read/write of intermediates, and distributes inputs/weights across the mesh via
  NoC row/column master-node multicast. The project's fixed 0.9875 correctness
  floor comes from the same paper's bf16-vs-fp32 context.}

\begin{figure}[t]
  \centering
  \begin{tikzpicture}[
      font=\small,
      kern/.style={draw=tensixnavy, very thick, fill=tensixlight,
                   rounded corners=1.5pt, minimum height=0.85cm,
                   minimum width=2.05cm, align=center, font=\footnotesize},
      dram/.style={draw=tensixslate, very thick, fill=tensixslate!12,
                   minimum height=0.75cm, minimum width=4.9cm, align=center,
                   font=\footnotesize\bfseries, text=tensixnavy},
      io/.style={-Stealth, thick, tensixnavy},
      saved/.style={-Stealth, thick, tensixamber, densely dashed},
      iolab/.style={font=\tiny, text=tensixnavy, align=left},
      paneltit/.style={font=\small\bfseries, text=tensixnavy},
    ]
  % ===================== UNFUSED =====================
  \node[paneltit] at (2.45, 3.35) {Unfused: 2 kernels, intermediate round-trips DRAM};
  \node[kern] (rn) at (1.25, 2.2) {RMSNorm};
  \node[kern] (mm) at (3.65, 2.2) {MatMul};
  \node[dram] (d1) at (2.45, 0.2) {DRAM};
  \draw[io] ($(d1.north)+(-1.9,0)$) -- ($(rn.south)+(-0.45,0)$)
    node[iolab, midway, left=1pt] {read $x$\\$16\,777\,216$\,B};
  \draw[io] ($(rn.south)+(0.45,0)$) -- ($(d1.north)+(-0.5,0)$)
    node[iolab, midway, right=0pt] {write $r$\\$16\,777\,216$\,B};
  \draw[io] ($(d1.north)+(0.55,0)$) -- ($(mm.south)+(-0.45,0)$)
    node[iolab, midway, left=1pt] {read $r + W$\\$117\,440\,512$\,B};
  \draw[io] ($(mm.south)+(0.45,0)$) -- ($(d1.north)+(2.0,0)$)
    node[iolab, midway, right=1pt] {write $y$\\$50\,331\,648$\,B};
  \node[font=\tiny, text=tensixslate, align=center] at (2.45, -0.55)
    {reads $134\,217\,728$\,B, writes $67\,108\,864$\,B};

  % ===================== FUSED =====================
  \node[paneltit] at (9.85, 3.35) {Fused: 1 kernel, intermediate stays in L1};
  \node[kern, minimum width=3.2cm] (fk) at (9.85, 2.2)
    {fused\_rmsnorm\_matmul\\[-1pt]{\tiny($r$ kept on-chip)}};
  \node[dram] (d2) at (9.85, 0.2) {DRAM};
  \draw[io] ($(d2.north)+(-1.0,0)$) -- ($(fk.south)+(-0.8,0)$)
    node[iolab, midway, left=1pt] {read $x + W$ once\\$117\,440\,512$\,B};
  \draw[io] ($(fk.south)+(0.8,0)$) -- ($(d2.north)+(1.0,0)$)
    node[iolab, midway, right=1pt] {write $y$ only\\$50\,331\,648$\,B};
  \node[font=\tiny, text=tensixslate, align=center] at (9.85, -0.55)
    {reads $117\,440\,512$\,B, writes $50\,331\,648$\,B};

  % elimination callout
  \node[draw=tensixamber, very thick, rounded corners=1.5pt, fill=white,
        font=\tiny, text=tensixnavy, align=center] at (9.85, 1.05)
    {eliminated round-trip: $r$ written + re-read\\ $= 16\,777\,216$\,B saved twice
     (reads \textbf{and} writes)};

  % separator
  \draw[tensixslate, thin, densely dotted] (6.3, -0.9) -- (6.3, 3.55);
  \end{tikzpicture}
  \caption{Unfused vs.\ fused dataflow for the running example. Unfused, the
    normalized intermediate $r$ ($16\,777\,216$\,B, same shape as the activation)
    is written to DRAM and re-read, adding one DRAM write and one DRAM read to the
    two-kernel execution. Fused, the intermediate stays in L1 between the stages:
    DRAM reads drop $134\,217\,728 \to 117\,440\,512$\,B and writes
    $67\,108\,864 \to 50\,331\,648$\,B --- the
    \texttt{estimated\_dram\_reads\_saved} $= 16\,777\,216$ figure recorded in
    \texttt{reports/rmsnorm\_matmul\_plan.json}. Byte figures are static analytic
    estimates from \texttt{docs/performance\_model.md}, not hardware measurements.}
  \label{fig:theory-fusion-dataflow}
\end{figure}

\subsection{NoC-aware sharding}

\texttt{-{}-tensixflow-plan-noc-sharding} maps the tensor shapes onto the 8x8 shard
grid and picks, tensor by tensor, how the data moves:

\begin{itemize}[nosep]
  \item \textbf{Activation: row-major sharding.} The leading (batch/M) dimension is
        row-major sharded across the grid's Y-cores, so each core row owns a
        contiguous band of rows and streams them tile by tile. The plan records
        this as \texttt{activation\_sharding: "row-major"}.
  \item \textbf{Weight: replicate or column-shard + multicast.} The weight planner
        compares the weight's footprint against the per-core circular-buffer budget
        of Figure~\ref{fig:theory-core-memory}. A weight that fits is
        \emph{replicated}: every core holds a full copy and multicast traffic is
        zero. A weight that does not fit is \emph{column-sharded} across the grid's
        X-cores and \emph{multicast} across the grid's rows via row master nodes,
        so it is moved once over the NoC rather than re-read from DRAM per core.
        The running example's weight ($[4096, 12288]$ bf16 $= 100\,663\,296$\,B
        $\gg$ 960\,KiB) takes the multicast path; the plan records
        \texttt{weight\_strategy: "multicast"}.
\end{itemize}

On top of the placement, the NoC traffic estimator produces a first-order static
account of both the unfused and the fused execution. For the running example it
reports unfused DRAM reads of
$16\,777\,216 + 16\,777\,216 + 100\,663\,296 = 134\,217\,728$ bytes
(activation $+$ intermediate $+$ weight) and unfused writes of
$16\,777\,216 + 50\,331\,648 = 67\,108\,864$ bytes (intermediate $+$ matmul
output), versus fused DRAM
reads $117\,440\,512$ bytes (activation + weight only, with the weight's
$100\,663\,296$ bytes counted once as NoC multicast) and writes $50\,331\,648$,
with the $16\,777\,216$-byte intermediate held L1-resident. These estimates are
also what the static performance model turns into the memory-vs-compute bound
classification; every figure is a shape-derived analytic estimate, never a measured
hardware number.

\subsection{Numerical validation theory}

The primary correctness metric is the Pearson correlation coefficient between the
simulator output vector $x$ and the golden output vector $y$:
\begin{equation}
  \mathrm{PCC}(x, y) \;=\;
  \frac{\sum_i (x_i - \bar{x})(y_i - \bar{y})}
       {\sqrt{\sum_i (x_i - \bar{x})^2}\;\sqrt{\sum_i (y_i - \bar{y})^2}},
  \label{eq:theory-pcc}
\end{equation}
with a \textbf{fixed correctness floor of 0.9875}: any run of either example graph
below the floor is a failure, and the floor is never lowered.%
\footnote{The floor and its context come from arXiv:2606.09879 (hardware bf16
  execution versus an fp32 reference, PCC kept above 98.75\%), the same grounding
  reference cited in Section~\ref{sec:theory}'s fusion subsection.}

The comparison is deliberately \emph{unequal-precision}: the simulator models the
bf16 spatial machine while the golden is fp32-exact, so Equation~\eqref{eq:theory-pcc}
measures a real quantity rather than an identity. Two effects of compiled execution
show up in the gap, and only in the gap:

\begin{enumerate}[nosep]
  \item \textbf{Precision cost.} The spatial machine rounds inputs to bf16, rounds
        every output tile back to bf16, and --- inside a fused kernel --- stores the
        intermediates between fused stages as bf16 in circular buffers
        (Figure~\ref{fig:theory-precision-flow}). Each round is a loss the
        fp32-exact golden never pays.
  \item \textbf{Tiled-execution cost.} Tile-by-tile accumulation adds partial sums
        in a different order than the composed whole-tensor evaluation of the
        golden, changing the rounding history even at fixed precision.
\end{enumerate}

If the golden also rounded its intermediates to bf16, the two paths would coincide
and PCC would collapse to 1.0, making any floor pass for nothing. The project
therefore keeps the golden fp32-exact and guards the bar with a layered
non-vacuity defense: a structural assertion that the golden's fused intermediates
are fp32 while the simulator output is bf16-representable, a bit-difference check
($\max_i |x_i - y_i| > 0$, i.e.\ PCC is strictly below 1.0), and an empirically
calibrated upper trip-wire above the measured good-case PCC. On the documented
seeds the measured minima are $\mathrm{PCC} = 0.99999316$ for
\texttt{rmsnorm\_matmul} and $\mathrm{PCC} = 0.99998885$ for \texttt{mlp\_block}
--- far above the 0.9875 floor, and strictly below 1.0, which is exactly the band
the rounding contract predicts.

\begin{figure}[t]
  \centering
  \begin{tikzpicture}[
      font=\footnotesize,
      simbox/.style={draw=tensixteal, very thick, fill=tensixlight,
                     rounded corners=1.5pt, align=center, font=\scriptsize,
                     text width=2.35cm, minimum height=1.15cm},
      goldbox/.style={draw=tensixnavy, very thick, fill=tensixlight,
                      rounded corners=1.5pt, align=center, font=\scriptsize,
                      text width=2.35cm, minimum height=1.15cm},
      flow/.style={-Stealth, thick, tensixslate},
      lanelab/.style={font=\footnotesize\bfseries, align=right},
    ]
  % ---- Simulator lane (bf16 spatial model) ----
  \node[lanelab, text=tensixteal] at (-0.4, 2.0) {Simulator\\(bf16 spatial\\model)};
  \node[simbox] (s1) at (1.7, 2.0) {inputs rounded\\to \textbf{bf16}};
  \node[simbox] (s2) at (4.5, 2.0) {tile compute:\\\textbf{fp32 accumulate}\\within each $32{\times}32$ tile};
  \node[simbox] (s3) at (7.3, 2.0) {each output tile\\rounded back\\to \textbf{bf16}};
  \node[simbox] (s4) at (10.1, 2.0) {fused-stage intermediates stored \textbf{bf16} in circular buffers};
  \node[simbox] (s5) at (12.9, 2.0) {output tiles\\\textbf{bf16}};
  \draw[flow] (s1) -- (s2); \draw[flow] (s2) -- (s3);
  \draw[flow] (s3) -- (s4); \draw[flow] (s4) -- (s5);

  % ---- Golden lane (fp32-exact) ----
  \node[lanelab, text=tensixnavy] at (-0.4, 0.0) {Golden\\(fp32-exact\\reference)};
  \node[goldbox] (g1) at (1.7, 0.0) {inputs\\\textbf{fp32}};
  \node[goldbox, text width=5.15cm] (g2) at (5.9, 0.0)
    {composed graph evaluated in \textbf{fp32}: fp32 accumulate, fp32 intermediates,
     \textbf{no} intermediate bf16 rounding};
  \node[goldbox] (g3) at (12.9, 0.0) {output\\\textbf{fp32}};
  \draw[flow] (g1) -- (g2); \draw[flow] (g2) -- (g3);

  % ---- Comparison ----
  \node[draw=tensixamber, very thick, rounded corners=1.5pt, fill=white,
        font=\footnotesize, text=tensixnavy, align=center] (cmp) at (7.3, -1.7)
    {$\mathrm{PCC}(\text{sim}_{\text{bf16 tiled}},\; \text{golden}_{\text{fp32 exact}})
      \;\geq\; 0.9875$\\[-1pt]{\scriptsize fixed floor; measured minima
      $0.99999316$ / $0.99998885$}};
  \draw[-Stealth, thick, tensixteal] (s5.south) -- ++(0,-0.6) -| (cmp.north east);
  \draw[-Stealth, thick, tensixnavy] (g3.south) -- ++(0,-0.45) -| (cmp.south west);
  \end{tikzpicture}
  \caption{Precision flow under the rounding contract. The simulator lane pays the
    bf16 spatial cost --- input rounding, fp32 accumulation within a tile, per-tile
    bf16 rounding, and bf16 circular-buffer intermediates between fused stages ---
    while the golden runs the composed graph fp32-exact throughout. The PCC between
    the two outputs therefore measures the real numerical cost of fusion plus tiled
    execution, rather than collapsing to a vacuous 1.0.}
  \label{fig:theory-precision-flow}
\end{figure}
